\section{Discussion}\label{sec:discussion}

The results support an evidence-guided approach to molecular property modelling, rather than a universal ranking of model families. Across five independent endpoints, descriptor-based, graph-based, and hybrid predictors each had tasks on which they performed well. Adding the 210 RDKit descriptors changed performance in the same direction in the framework and Chemprop implementations on four of the five tasks, and by a similar amount: a clear gain on FreeSolv and a change of a few percent or less elsewhere. On ESOL the two disagreed, but only because one Chemprop run diverged: its hybrid was better on nine of ten splits. This cross-implementation agreement, and the single visible failure, support treating descriptor utility as an empirical property of a task and protocol, and as a reason to keep per-split records rather than only summaries.

It is worth stating what this means for the chemist who has one dataset in hand. The default ``use the graph'' or ``use all descriptors'' habit is not wrong, but it is a guess, and the cost of testing it is now small. On phenacetin's endpoint the same graph encoder gave a test RMSE of $0.681$ without descriptors and $0.595$ with them on that split, a difference of $0.086$; averaged over ten splits the graph-only and full-descriptor means are within $0.001$ RMSE of each other on ESOL ($0.623$ versus $0.623$), which is exactly why a single split should not be used as a decision rule. The same holds for compact routes: the recommender gave the lowest mean on Lipophilicity ($0.581$ versus $0.591$ for random subsets), although the paired difference did not survive correction, and it did not separate from the control on any other endpoint. None of these numbers is a universal rule, and the framework deliberately does not turn any of them into one. What it does is make the alternatives cheap enough to compare on the user's own data, with the same split and the same record.

The first contribution is the coupling of feature evidence to model design, with a deliberate distinction between documentation and recommendation. The scorecard is an evidence-documentation and diagnostic layer: it records feature provenance, training support, target association, redundancy, and post-hoc range coverage, while the dataset audit makes duplicate structures and conflicting labels visible. The recommender is a separate candidate-generation heuristic that applies a fixed relevance and redundancy rule; it does not consume categorical scorecard evidence or automatically determine model roles. The recommender was reasonably stable across splits and often improved on random descriptor subsets when a Random Forest was trained directly on the selected descriptors. However, it did not outperform all selectors: Random Forest importance and mutual information ranked higher on average, and the recommender was not reliably superior to random subsets inside graph hybrids. We therefore interpret it as a transparent decision aid, not a new state-of-the-art selector or a measure of universal chemical importance. The scorecard's audit value does not depend on the recommender being the best-performing selector.

The second contribution is a configuration-driven modular architecture. Named feature blocks can be separated from routes and assigned to direct, fixed-prior, or learned-residual roles without changing feature definitions. The role ablation showed why this separation matters: direct fusion had the best mean on three of five tasks and the descriptor-only role on one, while graph only was best on BACE and within $0.001$ RMSE of the best route on Lipophilicity, so the ordering is endpoint-specific and often within between-split noise. The prior and residual design did not produce a statistically conclusive gain over graph only. Its value in this study is therefore architectural flexibility and a controlled way to test role assignments, rather than a claim that this role is consistently more accurate.

The third contribution is evaluation across molecular representations and different prediction endpoints. The same workflow accommodates fingerprints, computed descriptors, graph encoders, and hybrid inputs, and the benchmark spans scalar regression and binary classification. On the duplicate-controlled blood-brain barrier dataset, the framework and Chemprop hybrids had similar mean area under the receiver operating characteristic curve, while the Random Forest on RDKit descriptors had the highest mean. The exact ranking differed from the primary benchmark. Because the cleaned dataset and its split manifests also differed, this sensitivity analysis does not isolate duplicate leakage as the sole cause of the score change. Likewise, the benchmarks are separately trained single-endpoint experiments; they demonstrate cross-task workflow coverage, not the benefit of joint multi-task learning.

Several limitations remain. Ten split manifests still provide limited statistical power for differences of a few hundredths of a metric unit, and overlapping training sets make paired intervals descriptive rather than definitive. The primary blood-brain barrier benchmark retained duplicate structures, including groups with conflicting labels, and the audit found that tautomeric representations can produce different scaffold strings for the same InChIKey. The duplicate-controlled analysis addresses sensitivity to curation, but changes to data composition and splits prevent a causal attribution to leakage alone. Hyperparameters were fixed rather than exhaustively tuned. The descriptor-only reference and the framework's direct descriptor model were produced by different runners. The routing experiment now uses a harmonised checkpoint-selection rule across all five endpoints, with the complete checkpoint-level audit available in the repository. The ten-split summaries are descriptive and should not be interpreted as definitive evidence of small performance differences. Finally, the experiments use two-dimensional inputs and cover only regression and binary classification. Future work should evaluate additional datasets and prospective splits, strengthen duplicate-aware partitioning, test joint multi-task training where supported, and assess whether alternative feature evidence improves decisions beyond the present transparent heuristic.

